\subsection{Симулятор: что где находится}
\label{sec:simulator-tour}

Откройте \url{https://simulator.sakoryagin.ru}. При первом запуске появится короткая экскурсия --- её можно пройти или пропустить. После неё экран выглядит так:

\begin{shotcallouts}{01-scene.png}
    \callout{0.035}{0.955}{1}
    \callout{0.94}{0.90}{2}
    \callout{0.95}{0.73}{3}
    \callout{0.50}{0.19}{4}
    \callout{0.15}{0.07}{5}
    \callout{0.40}{0.07}{6}
    \callout{0.93}{0.064}{7}
\end{shotcallouts}

\begin{enumerate}[label=\bubble{\arabic*}]
    \item Кнопка панелей. За ней прячется «рельс» с панелями: \ui{Код}, \ui{Полёт} (телеметрия), \ui{Рой}, \ui{Сцена}, \ui{Учебник}, \ui{Пульт}, \ui{Настройки}.
    \item Компас осей. Ось \textbf{X} красная, \textbf{Y} зелёная, \textbf{Z} синяя и смотрит вверх. Щелчок по оси поворачивает камеру вдоль неё.
    \item Окна камер дрона: передняя и нижняя камеры, их же «видит» программа.
    \item \ui{Открыть редактор кода} --- сюда пишется программа. Позже эта кнопка исчезнет, редактор открывается панелью \ui{Код}.
    \item Камера: \ui{Земля} (вид сверху), \ui{Дрон} (камера следует за дроном), \ui{FPV} (от лица дрона), \ui{Свободно} (орбита вокруг точки или свободный полёт камеры).
    \item \ui{Объекты} и \ui{Окружение}: расстановка домов, маркеров, площадок; время суток и погода.
    \item \ui{Запустить}, \ui{Стоп} и \ui{Сбросить}. Сброс возвращает дрон на стартовую площадку и очищает след полёта.
\end{enumerate}

Камеру вращают левой кнопкой мыши, сдвигают правой, приближают колёсиком.

\subsection{Редактор кода}

Нажмите \ui{Открыть редактор кода}. Редактор занимает экран целиком, сцена при этом продолжает жить: программа работает и после закрытия редактора.

\begin{shotcallouts}{02-editor-lua.png}
    \callout{0.035}{0.83}{1}
    \callout{0.098}{0.922}{2}
    \callout{0.245}{0.922}{3}
    \callout{0.46}{0.922}{4}
    \callout{0.83}{0.922}{5}
    \callout{0.876}{0.922}{6}
    \callout{0.903}{0.922}{7}
    \callout{0.931}{0.922}{8}
    \callout{0.968}{0.922}{9}
\end{shotcallouts}

\begin{enumerate}[label=\bubble{\arabic*}]
    \item Рельс панелей: \ui{Код} (этот редактор), \ui{Полёт}, \ui{Рой}, \ui{Сцена}, \ui{Учебник}; внизу --- тема, \ui{Пульт}, \ui{Настройки}.
    \item Язык программы: \ui{Lua}, \ui{Python} или \ui{Blockly}. Переключение не стирает код: у каждого языка своя вкладка.
    \item Готовые примеры --- в том числе все программы из этой книги (папки \texttt{04\_language}, \texttt{05\_leds}, \texttt{06\_missions}).
    \item \ui{Открыть файл} с компьютера и \ui{Скачать код} в файл.
    \item Состояние дрона: \ui{На земле}, \ui{Взлёт}, \ui{В полёте}, \ui{Посадка}.
    \item Справочник API: описание и пример для каждой функции.
    \item Терминал: сюда попадает вывод \texttt{print} и сообщения автопилота.
    \item Уведомления: перед запуском симулятор проверяет программу и объясняет найденные ошибки здесь.
    \item Закрыть редактор и вернуться к сцене. Кнопки \ui{Запустить}, \ui{Стоп}, \ui{Сбросить} --- внизу справа, как и на сцене.
\end{enumerate}

\begin{hintbox}[Пять шагов до запуска]
Открыть редактор $\to$ выбрать язык $\to$ вставить или набрать программу $\to$ \ui{Запустить} $\to$ закрыть редактор и смотреть на сцену. Если дрон ведёт себя не так, как ожидалось, --- \ui{Стоп}, \ui{Сбросить}, и в терминале читать, что произошло.
\end{hintbox}
